\documentclass[11pt]{article}
\usepackage{setspace}
\usepackage{geometry}
\geometry{margin=0.85in}
\setstretch{1.0}

\begin{document}

\begin{flushright}
Zhongli Wang \\
+49 160 9541 6015 \\
nkwangzhongli@gmail.com \\
Bielefeld, Germany \\
\end{flushright}

\vspace{0.1cm}

\noindent
Recruiting Team \\
McKinsey \& Company \\
D\"usseldorf, Germany \\

\vspace{0.1cm}

\noindent
Application for Knowledge Analyst -- Energy Solutions \\

\vspace{0.05cm}

\noindent Dear Recruiting Team,

I am currently completing my Ph.D.\ in Economics at Bielefeld University and Universit\`a Cattolica del Sacro Cuore, with planned completion in September 2026. My doctoral work sits at the intersection of computational modeling, applied data analytics, simulation, and strategic decision-making under uncertainty. I am eager to apply this background to the energy sector, where market dynamics, operational constraints, decarbonization, and investment decisions are closely interconnected.

Across my doctoral research, I have developed end-to-end analytical workflows for complex economic and business systems. This includes data collection, data cleaning, feature engineering, model construction, simulation, validation, robustness checks, and the synthesis of quantitative findings into structured reports and presentations. In one dissertation project, I estimate an agent-based computational model using financial and operational data from AT\&T and T-Mobile, applying the Simulated Method of Moments and neural network-based surrogate modeling. This experience has strengthened my ability to work with complex datasets, build analytical tools, evaluate assumptions, and translate model outputs into decision-relevant insights.

My previous has focused strongly on the types of questions that are central to energy analytics: how firms make strategic investment decisions, how markets evolve under uncertainty, how operational capacity interacts with demand, and how quantitative models can support better decisions. I am particularly motivated by the opportunity to deepen my expertise in oil and gas operations, energy markets, and the energy transition while contributing my strengths in economic modeling, Python-based analytics, and structured problem solving. Topics such as decarbonization, operational performance, and risk management are especially appealing to me because they require both rigorous quantitative analysis and practical business judgment.
My research has trained me to develop reproducible analytical workflows, communicate complex findings clearly, and produce materials for both technical and non-technical audiences. I am comfortable working with statistical analysis, simulation-based methods, and data visualization, and I am motivated to further strengthen my Excel, Tableau, and energy-specific modeling skills.

I am attracted to McKinsey because of its distinctive combination of analytical rigor, client impact, apprenticeship culture, and global expertise. The Energy Solutions team is especially compelling to me because it transforms data, tools, benchmarks, and market perspectives into actionable insights for organizations across the energy value chain. I would be excited to contribute to the development of models, summary reports, benchmarks, and standalone presentations that help consultants and clients make better strategic, tactical, and operational decisions.

Thank you so much for considering my application! I would welcome the opportunity to discuss how my background in economics, quantitative modeling, and data-driven analysis can contribute to McKinsey's Energy Solutions team.

\vspace{0.15cm}

\noindent Sincerely, \\[0.5cm]
\noindent Zhongli Wang

\end{document}
